% Standalone problem document, generated from the adjacent README.md.
% Compile with XeLaTeX. Mathematical content is maintained in Markdown.
\documentclass[11pt,a4paper]{article}
\usepackage[fontsize=10.5pt]{scrextend}
\usepackage[margin=22mm,headheight=15pt,headsep=9mm,footskip=11mm]{geometry}
\usepackage{fontspec}
\setmainfont{texgyrepagella}[Extension=.otf,UprightFont=*-regular,BoldFont=*-bold,ItalicFont=*-italic,BoldItalicFont=*-bolditalic]
\setsansfont{texgyreheros}[Extension=.otf,UprightFont=*-regular,BoldFont=*-bold,ItalicFont=*-italic,BoldItalicFont=*-bolditalic]
\setmonofont{texgyrecursor}[Extension=.otf,UprightFont=*-regular,BoldFont=*-bold,ItalicFont=*-italic,BoldItalicFont=*-bolditalic,Scale=MatchLowercase]
\usepackage{amsmath,amssymb}
\usepackage{unicode-math}
\setmathfont{texgyrepagella-math.otf}
\usepackage{xcolor,microtype,enumitem,needspace,fancyhdr}
\usepackage{bookmark}
\definecolor{ink}{HTML}{17344B}
\definecolor{accent}{HTML}{197D80}
\definecolor{muted}{HTML}{596773}
\hypersetup{colorlinks=true,linkcolor=accent,urlcolor=accent,pdftitle={MF-05 - Local
Hölder continuity of the joint spectral
radius},pdfauthor={Open Problems in Numerical Linear Algebra}}
\urlstyle{same}
\setlength{\parindent}{0pt}
\setlength{\parskip}{5pt plus 1pt minus 1pt}
\linespread{1.04}
\setlength{\emergencystretch}{3em}
\widowpenalty=10000
\clubpenalty=10000
\displaywidowpenalty=10000
\allowdisplaybreaks[1]
\setlist{nosep,leftmargin=1.5em}
\providecommand{\tightlist}{\setlength{\itemsep}{0pt}\setlength{\parskip}{0pt}}
\providecommand{\pandocbounded}[1]{#1}
\pagestyle{fancy}
\fancyhf{}
\fancyhead[L]{\sffamily\footnotesize\color{muted}OPEN PROBLEMS IN NUMERICAL LINEAR ALGEBRA}
\fancyhead[R]{\sffamily\bfseries\footnotesize\color{accent}MF-05}
\fancyfoot[L]{\parbox[b]{0.9\textwidth}{\sffamily\fontsize{8}{10}\selectfont\color{muted}Editorial ratings; a bounded search cannot certify that a problem remains unsolved.\\Verification check: 2026-09-17}}
\fancyfoot[R]{\sffamily\footnotesize\color{muted}\thepage}
\renewcommand{\headrulewidth}{0.3pt}
\renewcommand{\headrule}{\hbox to\headwidth{\color{accent}\leaders\hrule height \headrulewidth\hfill}}
\makeatletter
\renewcommand{\section}{\Needspace{5\baselineskip}\@startsection{section}{1}{0pt}{12pt plus 2pt minus 2pt}{4pt}{\sffamily\bfseries\large\color{ink}}}
\renewcommand{\subsection}{\Needspace{4\baselineskip}\@startsection{subsection}{2}{0pt}{9pt plus 1pt minus 1pt}{3pt}{\sffamily\bfseries\normalsize\color{ink}}}
\makeatother
\setcounter{secnumdepth}{0}
\begin{document}
{\sffamily\small\color{muted}Matrix functions and stability\par}
\vspace{3pt}
{\raggedright\hyphenpenalty=10000\exhyphenpenalty=10000\sffamily\bfseries\fontsize{21}{24}\selectfont\color{ink}Local
Hölder continuity of the joint spectral radius\par}
\vspace{7pt}
{\sffamily\small\textbf{Difficulty:} challenging\quad\textbf{Importance:} interesting
to the community\par}
{\sffamily\small\bfseries\color{accent}Status: Lean verified\par}
\vspace{4pt}
{\color{accent}\hrule height 0.8pt}
\vspace{6pt}
\textbf{Rating rationale:} Challenging because reducible families
obstruct uniform perturbation estimates; community impact comes from
conditioning and robustness of stability computations.

\subsection{Resolution --- 2026-09-11}\label{resolution-2026-09-11}

\textbf{Affirmative resolution.} Matthew J. Colbrook's
\href{https://github.com/ajt60gaibb/OpenProblemsInNLA/blob/main/references/colbrook-jsr-growth-2026-09-11/manuscripts/uniform_growth_and_holder.pdf}{complete
manuscript, Theorem 2 and Corollary 7} proves the stronger uniform
estimate

\[|\widehat\rho(\mathcal M)-\widehat\rho(\mathcal N)|
\le d(2d+1)L^{1-1/d}d_H(\mathcal M,\mathcal N)^{1/d}\]

for any two nonempty compact real or complex matrix families in the
spectral-norm ball of radius \(L>0\). Choosing a common norm ball around
a fixed family gives the exact local two-family assertion below,
including reducible families and zero joint spectral radius. For \(d=1\)
the Lipschitz constant is one. The exponent \(1/d\) is sharp in general;
no Lipschitz lower-bound resolution of MF-06 is asserted.

The complete original proof passed
\href{https://github.com/ajt60gaibb/OpenProblemsInNLA/blob/main/references/colbrook-jsr-growth-2026-09-11/verification/reviews/MF-05-MF-07-review.md}{independent
Codex-agent review}.
\href{https://github.com/ajt60gaibb/OpenProblemsInNLA/blob/main/references/colbrook-jsr-growth-2026-09-11/manuscripts/uniform_growth_and_holder.tex}{Authored
TeX} ·
\href{https://github.com/ajt60gaibb/OpenProblemsInNLA/blob/main/references/colbrook-jsr-growth-2026-09-11/README.md}{Submission,
authorship and verification record}. The proof was developed with AI
assistance. The 11 September audit was informal; the subsequent Lean
verification is recorded below. No external human peer review is
claimed. The original statement and prior evidence below are retained,
and the ratings above are historical. This entry no longer contributes
to the open count.

\subsection{Lean proof and verification
evidence}\label{lean-proof-and-verification-evidence}

\textbf{The complete original target is Lean verified, 17 September
2026.} The
\href{https://github.com/sgstepaniants/OpenProblemsInNLA/blob/06b8cf49740205c4b7b0b71ee5c636855fbe26a6/matrix-functions-and-stability/MF-05/lean/Solution.lean}{immutable
proof} at revision \textit{06b8cf49740205c4b7b0b71ee5c636855fbe26a6}
passed
\href{https://github.com/sgstepaniants/OpenProblemsInNLA/actions/runs/35252365986/job/105307710518}{non-root
Linux run 35252365986}. The
\href{https://github.com/ajt60gaibb/OpenProblemsInNLA/blob/main/matrix-functions-and-stability/MF-05/lean/verification/linux-2026-09-17/README.md}{retained
execution evidence} authenticates all 193 submitted project inputs and
14 exported statements. Comparator statement matching, default-kernel
replay, standard transitive axioms, sandbox checks and rejection
controls all passed.

The main declaration \textit{NLA.MF05.canonical\_local\_holder} chooses
positive radius and constant before both compact complex matrix families
vary. It covers every positive dimension, infinite generating families,
reducibility, zero joint spectral radius and zero Hausdorff distance.
\textit{NLA.MF05.spectral\_hausdorff\_semantics} identifies the literal
maximum-of-suprema-of-infima metric, and
\textit{NLA.MF05.general\_root\_limit\_semantics} proves the full
chronological-product root-limit meaning of the radius, including radius
zero. \textit{NLA.MF05.uniform\_holder\_estimate} proves the displayed
uniform bound. The
\href{https://github.com/ajt60gaibb/OpenProblemsInNLA/blob/main/matrix-functions-and-stability/MF-05/lean/Challenge.lean}{14
exact contracts} and
\href{https://github.com/ajt60gaibb/OpenProblemsInNLA/blob/main/matrix-functions-and-stability/MF-05/lean/SourceCorrespondence.md}{source
correspondence} record the comparison with the unchanged original
target. Sharpness examples and the better scalar Lipschitz constant are
not additional formal claims.

Formalization and verification submission: \textbf{George Stepaniants},
Department of Computing and Mathematical Sciences, California Institute
of Technology. \textbf{Matthew J. Colbrook}, Department of Applied
Mathematics and Theoretical Physics, University of Cambridge, retains
mathematical proof credit. Two independent statement reviews preceded
implementation; two independent nonauthor final source reviews approved
the complete proof.
\href{https://github.com/ajt60gaibb/OpenProblemsInNLA/blob/main/matrix-functions-and-stability/MF-05/lean/reviews/README.md}{Review
records} disclose substantial OpenAI Codex assistance and the reviewers'
exact scopes.

The project pins Lean 4.33.1,
\href{https://github.com/leanprover-community/mathlib4/tree/0df444a360eaa60ab8c11dca51a86af692955474}{Mathlib}
and
\href{https://github.com/alerad/leancert/tree/621a43d7cf21f87872392a01e874f2f1dbddc926}{LeanCert}.
The consumed LeanCert certificate proves \(0<1/2<1\) in kernel mode; the
remaining argument is symbolic. Every export has only \textit{propext},
\textit{Classical.choice} and \textit{Quot.sound} as transitive axioms.
Run \textit{lake build} in
\textit{matrix-functions-and-stability/MF-05/lean};
\href{https://github.com/ajt60gaibb/OpenProblemsInNLA/blob/main/matrix-functions-and-stability/MF-05/lean/README.md}{full
reproduction instructions} distinguish the actual local builds from the
final Linux check.

\subsection{Context and notation}\label{context-and-notation}

Let \(\mathcal H_d\) denote the nonempty compact subsets of
\(\mathbb C^{d\times d}\). Use the spectral norm and its Hausdorff
distance

\[d_H(\mathcal M,\mathcal N)=\max\left\{
\sup_{A\in\mathcal M}\inf_{B\in\mathcal N}\|A-B\|_2,
\sup_{B\in\mathcal N}\inf_{A\in\mathcal M}\|A-B\|_2\right\}.\]

The joint spectral radius is

\[\widehat\rho(\mathcal M)=\lim_{k\to\infty}
\max_{A_1,\ldots,A_k\in\mathcal M}\|A_k\cdots A_1\|_2^{1/k}.\]

These definitions also apply to finite real matrix sets. The ordinary
spectral radius of one matrix is written \(\rho(A)\).

\subsection{Problem statement}\label{problem-statement}

For every \(d\ge1\) and \(\mathcal M_0\in\mathcal H_d\), do there exist
\(r,C>0\) such that

\[|\widehat\rho(\mathcal M)-\widehat\rho(\mathcal N)|
\le C d_H(\mathcal M,\mathcal N)^{1/d}\]

whenever \(d_H(\mathcal M,\mathcal M_0)< r\) and
\(d_H(\mathcal N,\mathcal M_0)< r\)?

\subsection{Reference and status
evidence}\label{reference-and-status-evidence}

Epperlein and Wirth, \href{https://arxiv.org/html/2311.18633v2}{The
joint spectral radius is pointwise Hölder continuous}, Linear Algebra
Appl. 704 (2025), 92--122,
\href{https://doi.org/10.1016/j.laa.2024.09.016}{DOI}, §2, Conjecture 3
(L1). Pointwise results in that paper do not prove this local assertion.

\subsection{Additional status
evidence}\label{additional-status-evidence}

Searches combining ``joint spectral radius'' with ``local Hölder'',
``Lipschitz lower'', ``trajectory bounds'', the authors' names, and
2025/2026 found no later resolution. These searches supplement the
explicit 2025 conjectures; they do not establish exhaustiveness. The
three entries are separately named assertions in the source, not a count
of dimensional cases.

\subsection{Audit --- 2026-09-10}\label{audit-2026-09-10}

Rechecked \href{https://arxiv.org/html/2311.18633v2}{Epperlein--Wirth,
§1 and Conjecture 3 (L1)}. Local Lipschitz continuity settles the
irreducible-family subcase, but the displayed exponent for arbitrary
compact families remains open. Author and local-Hölder follow-up
searches found no resolution; the weaker general pointwise results do
not suffice.
\end{document}
